% Auto-generated by ResearchClaw — NeurIPS 2025 template
% Style file: https://media.neurips.cc/Conferences/NeurIPS2025/Styles.zip
\documentclass{article}
\usepackage[preprint]{neurips_2025}
\usepackage{hyperref}
\usepackage{url}
\usepackage{booktabs}
\usepackage{amsfonts}
\usepackage{amsmath}
\usepackage{nicefrac}
\usepackage{microtype}
\usepackage{graphicx}
\usepackage{natbib}
\usepackage[utf8]{inputenc}
\usepackage[T1]{fontenc}

\title{An Empirical Study of Privacy-Utility Tradeoffs in Anonymized Error Trace Sharing}

\author{Anonymous}

\begin{document}
\maketitle

\begin{abstract}
Developer tools can be significantly improved by analyzing error patterns aggregated from a large user base. However, collecting raw error traces poses a substantial privacy risk. This paper empirically investigates the tradeoff between diagnostic utility and privacy when anonymizing software error traces. Existing approaches often apply uniform privacy policies that can unnecessarily degrade data utility. We propose and evaluate two distinct structure-aware anonymization techniques: pseudonymization via token hashing and generalization via frequency-based thresholding (a form of k-anonymity). We measure utility via a pairwise bug matching task (Matching Accuracy) and privacy via a simulated project-linkage attack (Leakage Risk) on a synthetic dataset with varying data heterogeneity. Our results show that in a challenging high-diversity environment, the generalization-based approach provided a superior balance, achieving a Matching Accuracy of 0.6492 and a Leakage Risk of 0.0847. In contrast, token hashing yielded a lower accuracy of 0.6075 with a significantly higher risk of 0.2385. These findings suggest that the optimal anonymization strategy depends on data characteristics, and that generalization can offer a more robust solution for heterogeneous data sources, providing actionable guidance for designing privacy-preserving telemetry systems.
\end{abstract}

\section{Introduction}

\label{sec:introduction}

Modern software development increasingly relies on collective intelligence systems to improve developer productivity and software quality. Services like Sentry, Bugsnag, and Windows Error Reporting aggregate error reports from millions of users, enabling developers to discover, prioritize, and debug issues that would be invisible at the scale of a single user or organization. The core value proposition of these systems is their ability to cluster vast streams of incoming error reports into groups of unique, actionable bugs. A developer can then see that a single underlying fault is affecting thousands of users, justifying the allocation of engineering resources to fix it.

However, the power of these systems is predicated on collecting rich, detailed diagnostic data, which creates a fundamental tension with user privacy. A typical software error trace, or stack trace, is a sequence of function calls leading to a fault. These traces often contain personally identifiable information (PII) or commercially sensitive data, such as user-specific file paths (\texttt{/home/alice/my\_secret\_project/main.py}), internal API endpoints, or data values printed in an error message. The unmitigated collection of such data represents a significant privacy liability and can erode user trust, hindering the adoption of these valuable tools. This creates a critical need for anonymization techniques that can provably protect user privacy while preserving the diagnostic utility of the collected data.

The central challenge lies in navigating the privacy-utility tradeoff. Existing approaches to anonymizing developer telemetry often treat the error trace as a flat string or a bag of words. They apply monolithic privacy techniques, such as adding noise to token embeddings to satisfy differential privacy \cite{zheng2025benchmarking}, or applying simple redaction rules. These methods are blunt instruments. They fail to recognize that different parts of an error trace have vastly different sensitivity and utility. For instance, a function name from a public library (\texttt{numpy.linalg.inv}) is low-sensitivity but high-utility for clustering, whereas a local variable name (\texttt{user\_password}) is high-sensitivity but low-utility. A uniform anonymization strategy either fails to protect sensitive data or needlessly degrades high-utility signals.

This paper argues for a more nuanced, structure-aware approach. We posit that by parsing an error trace into its semantic components (e.g., library code vs. user code, function names vs. variable names), we can apply targeted anonymization strategies that are better aligned with the specific privacy risk and utility contribution of each component. Rather than seeking a single "best" method, we conduct an empirical study to compare two distinct structure-aware strategies---pseudonymization and generalization---across different data environments. This leads to our primary research questions:

* \textbf{RQ1:} How does the utility of anonymized error traces, measured by their ability to support the identification of similar bugs, change under different anonymization strategies?
* \textbf{RQ2:} How does privacy, measured by the difficulty of re-identifying the source of a trace, change under these strategies?
* \textbf{RQ3:} How does the relative performance of different anonymization strategies (pseudonymization vs. generalization) change based on the heterogeneity of the underlying data source?

To answer these questions, we propose an evaluation methodology that compares two anonymization techniques: token hashing (a form of pseudonymization) and frequency-based generalization (a form of k-anonymity). Both utility (pairwise matching accuracy) and privacy (resistance to a linkage attack) are measured quantitatively, allowing us to plot their respective positions on the privacy-utility plane and directly compare the methods under different data conditions.

Our main contributions are:
* A comparative analysis of two structure-aware anonymization techniques for error traces: \textbf{token hashing} (pseudonymization) and \textbf{frequency-based generalization} (a k-anonymity variant).
* An evaluation methodology that jointly measures \textbf{utility}, via Matching Accuracy on a pairwise bug classification task, and \textbf{privacy}, via a simulated project-linkage attack, to enable a direct, quantitative comparison of the tradeoff.
* An empirical study on a synthetic dataset with controlled diversity that demonstrates how data heterogeneity impacts the performance of different anonymization strategies.
* A key finding from our study that \textbf{generalization offers a more robust privacy-utility balance than token hashing in high-diversity data environments}, suggesting its suitability for large-scale error aggregation platforms.

This paper is organized as follows. Section 2 reviews related work in data privacy and trace anonymization. Section 3 details our anonymization methods. Section 4 describes the experimental design, dataset, and evaluation metrics. Section 5 presents our quantitative results, and Section 6 discusses their implications. Section 7 outlines the study's limitations, and finally, Section 8 concludes with a summary and directions for future work.

\section{Related Work}

\label{sec:related_work}

Our research is situated at the intersection of privacy-preserving data publishing, software engineering, and machine learning. We structure our review of prior work into three corresponding areas.

\subsection{Privacy-Preserving Data Publishing and Anonymization}

\label{sec:privacy_preserving_data_publishing_and_a}

The study of privacy-utility tradeoffs is a well-established field. Early work focused on syntactic privacy models like \textit{k}-anonymity, which requires that each record in a released dataset be indistinguishable from at least \textit{k}-1 other records \cite{sakuma2017recommendation}. Subsequent extensions like \textit{l}-diversity and \textit{t}-closeness were proposed to mitigate attribute disclosure attacks. While foundational, these models are often difficult to apply to high-dimensional, semi-structured data like stack traces and do not always provide the rigorous, provable guarantees of modern techniques \cite{domingoferrer2025actual}. Our work operationalizes a form of \textit{k}-anonymity through generalization, evaluating its empirical effectiveness in a specific software engineering context rather than focusing on its formal properties.

Differential Privacy (DP) has emerged as the gold standard for statistical privacy \cite{desfontaines2019differential}. DP provides a formal guarantee that the outcome of a computation is insensitive to the presence or absence of any single individual's data. It is typically achieved by adding calibrated noise to query results or model parameters \cite{jiang2025balancing}. DP has been successfully applied in numerous domains, from federated learning to generating synthetic data. While DP provides the strongest theoretical guarantees, applying it monolithically can sometimes result in poor utility \cite{zhao2020many}, a problem our work explores through alternative, structure-aware methods. A comprehensive evaluation of pragmatic techniques like ours against a formal DP baseline remains a crucial direction for future work. Visualizing these tradeoffs is also an active area of research.

\subsection{Log and Trace Anonymization in Software Engineering}

\label{sec:log_and_trace_anonymization_in_software_}

Within software engineering, there is a history of work on anonymizing execution traces for security and debugging purposes. Early work focused on network traces, developing techniques like IP address truncation and prefix-preserving anonymization to balance privacy against utility for tasks like anomaly detection. These studies established the core tradeoff but often relied on ad-hoc utility metrics and did not offer formal privacy guarantees.

More recent work has begun to explore anonymization for software analytics data, such as telemetry from developer tools. Some approaches focus on pseudonymization \cite{yermilov2023privacy} or using large language models (LLMs) to redact sensitive information while preserving natural language context. For example, \cite{elkoumy2022libra} uses subsampling to anonymize event logs for process mining. Our work differs from these by focusing on a systematic, quantitative comparison of different \textit{types} of anonymization strategies (pseudonymization vs. generalization) and evaluating their performance as a function of data heterogeneity. We aim to provide guidance on \textit{which} type of strategy is more suitable for a given environment, a goal more aligned with foundational data privacy research \cite{avent2019automatic}.

\subsection{Measuring Privacy and Utility}

\label{sec:measuring_privacy_and_utility}

A key challenge in this domain is defining and measuring privacy and utility \cite{wagner2015technical}. Utility is task-dependent. For our task of collective error analysis, the primary utility lies in correctly grouping similar error reports. Therefore, we follow prior work in bug report analysis and use a classification task as a proxy for utility. Specifically, we measure the accuracy of a model trained to determine if two anonymized traces belong to the same ground-truth bug.

Measuring privacy is more complex, as it requires defining a threat model. A formal guarantee like DP is one approach, but it measures privacy in the abstract. An alternative, complementary approach is to measure privacy empirically by simulating a specific attack. Membership inference attacks, which aim to determine if a specific record was in the training data, are a common benchmark \cite{nasr2018comprehensive}. Linkage attacks, which aim to link an anonymized record to an external dataset, are also widely studied \cite{rastogi2006boundary}. Our privacy metric is based on a simulated linkage attack tailored to our domain: an adversary attempts to link an anonymized error trace back to the open-source project it originated from. This provides a concrete, interpretable measure of re-identification risk that complements theoretical guarantees. Other work has explored similar empirical evaluations for different data types. Our work is novel in its direct comparison of these metrics across different anonymization philosophies applied to structured error data.

\section{Method}

\label{sec:method}

In this section, we formalize the problem of error trace anonymization and detail the specific methods evaluated in our study. Our goal is to compare anonymization functions $\mathcal{A}$ that transform a raw trace $T$ into an anonymized trace $T' = \mathcal{A}(T)$ in order to understand their relative impact on diagnostic utility $U(T')$ and privacy risk $R(T')$.

\subsection{Problem Formulation}

\label{sec:problem_formulation}

We define a raw stack trace $T$ as a sequence of frames, $T = (f_1, f_2, \dots, f_L)$, where $L$ is the trace depth. Each frame $f_i$ is a tuple containing structured information, primarily the file path $p_i$ and the function name $m_i$. Thus, $f_i = (p_i, m_i)$. The core of our task is to design and compare transformations on the set of paths $\{p_i\}$ and methods $\{m_i\}$ that obscure sensitive information while preserving the structural patterns necessary for grouping related errors.

The overall problem is to explore the outcomes of different strategies $\mathcal{A}$ on the bi-objective optimization problem:
$$
\max\_{\mathcal{A}} \left( U(\{\mathcal{A}(T)\}\_{T \in \mathcal{D}}), -R(\{\mathcal{A}(T)\}\_{T \in \mathcal{D}}) \right)
$$
where $\mathcal{D}$ is a corpus of traces, $U$ is a utility metric, and $R$ is a privacy risk metric. We investigate this by evaluating discrete, representative points on the potential solution space.

\subsection{Anonymization Techniques}

\label{sec:anonymization_techniques}

We evaluate three distinct anonymization strategies: a non-private baseline and two structure-aware techniques based on pseudonymization and generalization, respectively.

\begin{enumerate}
 \item \textbf{No Anonymization (No-Op)}: This serves as a baseline representing the maximum possible utility and maximum privacy risk. The anonymization function is the identity function, $\mathcal{A}_{none}(T) = T$. All original file paths and function names are preserved. This method provides the upper bound on utility, as it retains all original information for the downstream task.
 \item \textbf{Structure-Preserving Token Hashing (Pseudonymization)}: This technique aims to break the link between identifiers and their real-world semantic meaning while preserving the ability to group identical traces. It is a form of pseudonymization. The anonymization function $\mathcal{A}_{hash}$ processes each token (file path segment or function name) $t$ within the trace. Each unique token is replaced by its salted hash: $t' = \text{SHA256}(t + \text{salt})$. A key feature of this method is that the salt is unique per project, but consistent for all traces originating from that same project. This allows for perfect matching of tokens \textit{within} a project's error reports but prevents trivial cross-project linkage based on identifier names. The overall structure of the trace (sequence of frames and tokens) is perfectly preserved. This method preserves the cardinality of the original token space.
 \item \textbf{Generalization via Frequency Thresholding (k-Anonymity Variant)}: This method seeks to provide privacy by ensuring that rare, and therefore potentially identifying, tokens are obscured. It implements a concept from \textit{k}-anonymity by grouping rare items into a single category. Our implementation, $\mathcal{A}_{k-anon}$, achieves this through rule-based generalization of identifier tokens. First, we perform a simple structural generalization on file paths, replacing common prefixes for system libraries or user home directories with generic tokens like \texttt{[SYS\_LIB]} and \texttt{[USER\_HOME]}. Second, and more importantly, we apply a frequency-based generalization to all tokens (function names, class names, remaining path segments). We build a frequency map of all tokens in a reference corpus (the training set). Any token that appears fewer than \textit{k} times is replaced with a generic placeholder token, e.g., \texttt{[RARE\_TOKEN]}. This process reduces the uniqueness of individual traces by reducing the cardinality of the token space, forcing them into larger equivalence classes. Unlike hashing, which is a one-to-one mapping, generalization is a many-to-one mapping, inherently increasing ambiguity to protect privacy.
\end{enumerate}

\subsection{Trace Representation}

\label{sec:trace_representation}

To evaluate utility and privacy, we must convert the anonymized, structured traces into a format suitable for machine learning. We first transform each trace $T'$ into a single document by concatenating all its tokens (anonymized file path segments and function names). We then use a standard Term Frequency-Inverse Document Frequency (TF-IDF) vectorizer to convert the entire corpus of trace documents into a matrix of numerical feature vectors. This approach captures the relative importance of different tokens for distinguishing between traces. The resulting vectors serve as input to the downstream tasks used for evaluating utility and privacy. To prevent data leakage, the TF-IDF vectorizer is fit only on the training portion of the dataset, and this fitted vectorizer is used to transform the test set.

\section{Experiments}

\label{sec:experiments}

We designed an empirical study to quantify and compare the privacy-utility tradeoffs of the proposed anonymization methods. The experiments are designed to answer our research questions by measuring the impact of each method on a utility task and a privacy-breaching task under different data conditions.

\subsection{Dataset and Regimes}

\label{sec:dataset_and_regimes}

To ensure a controlled environment for comparison, we constructed a synthetic dataset of 10,000 Python error traces. The dataset generation process was designed to be reproducible and to simulate two distinct real-world scenarios, or regimes, to test the robustness of our methods.

\textbf{Data Generation Process:}
\begin{enumerate}
 \item \textbf{Codebase Simulation:} We first defined a vocabulary of 1,000 unique function names, partitioned into "library" functions (common, shared across projects) and "application" functions (specific to a project). We simulated 50 distinct projects.
 \item \textbf{Bug Template Definition:} We defined 500 ground-truth bug "templates." Each template consists of a unique call stack sequence of a specific depth, composed of function names drawn from the vocabulary. Bugs could be project-specific (using only that project's application functions and library functions) or cross-project (involving calls between simulated projects).
 \item \textbf{Trace Generation:} For each of the 500 bug templates, we generated 20 individual trace instances, for a total of 10,000 traces. Each instance was created by taking the bug template's call stack and instantiating it with realistic-looking file paths. Paths included variations in user home directories (\texttt{/home/user1/}, \texttt{/home/user2/}) and installation paths (\texttt{/opt/venv/}, \texttt{/usr/lib/python3.9/}) to simulate real-world noise.
\end{enumerate}

This process yields a dataset with known ground-truth labels for both the underlying bug (for utility measurement) and the source project (for privacy measurement). We then organized this dataset into two regimes:

* \textbf{Low-Diversity Regime}: This regime simulates error reports from a single, large, monolithic application. All 10,000 traces were generated as if originating from a single "project," sharing a large common pool of application functions. The traces exhibit significant overlap in their library dependencies and internal code structure. This represents an easier case for clustering, as shared stack frames are common.
* \textbf{High-Diversity Regime}: This regime simulates error reports aggregated from hundreds of distinct, smaller projects. The 10,000 traces were generated across the 50 distinct simulated projects. Traces are more heterogeneous, with less overlap in their call stacks. This presents a more challenging scenario for both utility and privacy tasks.

For each regime, we applied the three anonymization methods described in Section 3.2: No Anonymization, Token Hashing, and Generalization. This results in a total of $2 \times 3 = 6$ experimental conditions.

\subsection{Evaluation Metrics}

\label{sec:evaluation_metrics}

We define two primary metrics to quantitatively assess the privacy-utility tradeoff.

* \textbf{Utility: Matching Accuracy}. Utility is defined as the ability to correctly identify whether two error traces belong to the same underlying bug. We measure this using a pairwise classification task. From the set of all traces, we create a balanced dataset of trace pairs, with an equal number of positive pairs (two traces from the same bug cluster) and negative pairs (two traces from different bug clusters). Given the TF-IDF vectors of two traces in a pair, we concatenate their vectors to form a single feature vector. A logistic regression classifier is trained on these feature vectors to predict if the pair shares the same ground-truth bug label (a binary prediction). \textbf{Matching Accuracy} is the accuracy of this classifier on a held-out test set. A higher accuracy (range [0, 1]) indicates that the anonymized representation retains more useful signal for grouping related errors. An accuracy of 1.0 means the representation perfectly separates all bug clusters.
* \textbf{Privacy: Leakage Risk}. Privacy is measured by quantifying the risk of re-identification. We simulate a linkage attack where an adversary's goal is to identify the source project of a given anonymized trace. We train a multi-class logistic regression classifier to predict a trace's source project (from the 50 distinct projects in the High-Diversity regime) based on its anonymized TF-IDF vector. \textbf{Leakage Risk} is the accuracy of this attacker model on a held-out test set. A lower value is better. A value of 1.0 indicates complete leakage (no privacy), while a value close to random chance (1/50 = 0.02 in our setup) indicates strong privacy protection.

\subsection{Experimental Setup}

\label{sec:experimental_setup}

All anonymization and evaluation procedures were implemented in Python using the Scikit-learn library. The dataset was split into an 80\% training set and a 20\% testing set. The TF-IDF vectorizer and all classifiers were trained only on the training set. The results reported are based on a single experimental run (n=1); consequently, we do not report confidence intervals or conduct statistical significance tests. The findings should be interpreted as preliminary observations from a pilot study.

\begin{table}[ht]
\centering
\begin{tabular}{l l l}
\toprule
\textbf{Parameter} & \textbf{Value} & \textbf{Description} \\
\midrule
\textbf{TF-IDF} & & \\
\texttt{max\_features} & 5000 & Maximum number of features in the vocabulary. \\
\texttt{min\_df} & 2 & Minimum document frequency for a token to be included. \\
\textbf{Generalization} & & \\
\texttt{k} & 5 & Minimum frequency for a token to be preserved. \\
\textbf{Token Hashing} & & \\
\texttt{hash\_algorithm} & SHA256 & Hashing algorithm for pseudonymization. \\
\texttt{salt} & Per-project & Salt is unique for each project, shared for traces within it. \\
\textbf{Classifiers} & & \\
\texttt{model} & Logistic Regression & Model for both Matching Accuracy and Leakage Risk tasks. \\
\texttt{solver} & \texttt{lbfgs} & Optimization algorithm for the logistic regression model. \\
\texttt{max\_iter} & 1000 & Maximum iterations for solver convergence. \\
\bottomrule
\end{tabular}
\caption{Table Hyperparameter settings for the experimental setup.}
\label{tab:1}
\end{table}

\section{Results}

\label{sec:results}

Our experiments provide a preliminary quantitative analysis of the privacy-utility tradeoff for the evaluated anonymization techniques. The results, based on a single experimental run, are presented across two distinct data regimes: a low-diversity setting simulating a monolithic application and a high-diversity setting simulating an aggregation of many small projects.

Table 2 summarizes the performance of each method in terms of Matching Accuracy (utility) and Leakage Risk (privacy). Matching Accuracy measures the ability to group related errors, with higher values being better. Leakage Risk measures the re-identifiability of traces, with lower values indicating better privacy.

\textbf{Table 2: Privacy-Utility Tradeoffs of Anonymization Methods.} Results from a single run. Higher Matching Accuracy is better for utility. Lower Leakage Risk is better for privacy. Bold indicates the best-performing \textit{anonymization method} for each metric within each regime.

\begin{table}[ht]
\centering
\begin{tabular}{l l l l}
\toprule
\textbf{Regime} & \textbf{Method} & \textbf{Matching Accuracy (Utility ↑)} & \textbf{Leakage Risk (Privacy ↓)} \\
\midrule
\textbf{Low-Diversity} & No Anonymization & 1.0000 & 1.0000 \\
 & Token Hashing & \textbf{0.8581} & 0.2379 \\
 & Generalization (k=5) & 0.8260 & \textbf{0.0531} \\
 & & & \\
\textbf{High-Diversity} & No Anonymization & 1.0000 & 1.0000 \\
 & Token Hashing & 0.6075 & 0.2385 \\
 & Generalization (k=5) & \textbf{0.6492} & \textbf{0.0847} \\
\bottomrule
\end{tabular}
\caption{Table 2}
\label{tab:2}
\end{table}

From these initial results, we draw several key observations:

\textbf{The Fundamental Tradeoff is Apparent:} The \texttt{No Anonymization} baseline confirms the core tension. It achieves perfect utility (Matching Accuracy of 1.0000) by retaining all information, but at the cost of complete privacy failure (Leakage Risk of 1.0000). This establishes the upper bound for utility and the lower bound for privacy that other methods are measured against.

\textbf{Anonymization is Effective at Reducing Risk:} Both \texttt{Token Hashing} and \texttt{Generalization} substantially reduce Leakage Risk compared to the baseline. The \texttt{Generalization} method appears particularly effective, reducing risk from 1.0000 to 0.0531 in the Low-Diversity regime and to 0.0847 in the High-Diversity regime. This demonstrates that structured generalization can be a powerful tool for privacy protection.

\textbf{Regime-Dependent Performance:} The effectiveness of each method is highly dependent on the data characteristics.

In the \textbf{Low-Diversity} regime, which simulates a single large application, the data is more homogeneous. In this run, \texttt{Token Hashing} maintained a higher utility than \texttt{Generalization} (Matching Accuracy of 0.8581 vs. 0.8260). However, this utility gain came with a significantly higher privacy risk (Leakage Risk of 0.2379 vs. 0.0531). In this context, \texttt{Generalization} appears to offer a more favorable tradeoff, achieving strong privacy with a modest utility cost.

In the \textbf{High-Diversity} regime, which simulates aggregating errors from many different projects, the data is more heterogeneous. This setting proved more challenging for utility preservation for both methods. Here, we observed an interesting reversal in relative utility. In this run, \texttt{Generalization} not only provided superior privacy compared to \texttt{Token Hashing} (Leakage Risk of 0.0847 vs. 0.2385) but also yielded higher utility (Matching Accuracy of 0.6492 vs. 0.6075). This suggests that when traces are inherently dissimilar, the generalization strategy may be more effective at preserving the core structural signals for matching than the pseudonymization approach of \texttt{Token Hashing}.

Overall, this pilot study suggests that \texttt{Generalization} is a robust strategy that provides strong privacy protection across different data environments. While \texttt{Token Hashing} can preserve more utility in homogeneous settings, \texttt{Generalization} appears to offer a more consistently balanced and often superior profile, especially in more realistic, high-diversity scenarios.

\section{Discussion}

\label{sec:discussion}

The results of our initial exploration provide valuable insights into the complex interplay between privacy and utility in collective error reporting. Our findings suggest that the choice of anonymization strategy is not merely a technical detail but a critical decision that should be informed by the characteristics of the underlying data.

Our primary observation is that structured anonymization techniques can drastically reduce re-identification risk while preserving a large fraction of diagnostic utility. In the Low-Diversity setting, our \texttt{Generalization} method reduced Leakage Risk to 0.0531 while retaining a Matching Accuracy of 0.8260. This is a promising result, suggesting that systems can move far away from the dangerous default of collecting raw data (1.0 risk, 1.0 utility) to a much safer operating point without rendering the data useless. This challenges the notion that any meaningful anonymization leads to a catastrophic loss of utility. Instead, our data points to a curve where significant privacy gains can be achieved for a manageable utility cost.

The differing performance of \texttt{Token Hashing} and \texttt{Generalization} across the two regimes is particularly illuminating. The superiority of \texttt{Generalization} in the challenging High-Diversity setting is a key finding of this study. We hypothesize that this is because \texttt{Token Hashing} (a form of pseudonymization) preserves the \textit{cardinality} of the identifier space. In a high-diversity corpus with many distinct projects, this results in a vast number of unique hashed tokens, making the TF-IDF vectors sparse and brittle, which may harm the performance of downstream models. In contrast, \texttt{Generalization} (a form of k-anonymity) actively \textit{reduces} the identifier cardinality by mapping all rare tokens to a single \texttt{[RARE\_TOKEN]} placeholder. This acts as a form of automated feature engineering, forcing the matching model to rely on more common, structural elements of the traces, which appears to be a more robust signal in heterogeneous data.

These findings have direct practical implications for developers of telemetry systems. They suggest that a one-size-fits-all approach is suboptimal. For platforms collecting data from a single product (e.g., a SaaS application), \texttt{Token Hashing} might be a viable option if the utility premium is critical and the higher leakage risk is acceptable within the system's threat model. However, for large-scale aggregators collecting data from thousands of open-source or commercial projects (e.g., a package manager's error reporting service), a generalization-based approach like our \texttt{Generalization} method appears to be a more robust and safer choice, offering both better utility and better privacy in our high-diversity experiment.

Our work aligns with prior research in privacy-preserving data publishing that highlights the fragility of syntactic privacy models \cite{domingoferrer2025actual}. However, by grounding our evaluation in a concrete, domain-specific task, we provide an empirical counterpoint to purely theoretical critiques. Our results suggest that simple, pragmatic techniques can offer an effective balance in real-world software engineering scenarios, motivating further research into more advanced and hybrid methods. It is also important to note that our privacy evaluation is limited to a single linkage attack scenario and does not account for other potential threats, such as attribute inference, which should be considered in a complete threat model.

\section{Limitations}

\label{sec:limitations}

While this study introduces a valuable framework and provides initial insights, its limitations must be acknowledged to guide future work and to properly contextualize our findings.

* \textbf{Single Experimental Run:} The empirical results are based on a single experimental run (n=1). Consequently, we cannot report confidence intervals or perform statistical significance tests. The observed differences in performance between methods should be considered preliminary and require validation across multiple runs with different random seeds to ensure their stability and generalizability.
* \textbf{Synthetic Dataset:} The study was conducted on a synthetically generated dataset. Although designed with controllable parameters to capture properties of low- and high-diversity environments, it may not fully represent the complexity, noise, and idiosyncrasies of real-world error traces, such as those arising from complex library versioning conflicts or multithreaded applications. Generalizing these findings requires further validation on large-scale, real-world corpora of error reports.
* \textbf{Simplistic Anonymization Methods:} The implemented anonymization techniques, while representative of their respective classes (pseudonymization and generalization), are simplistic. The \texttt{Generalization} method uses a global frequency threshold, which may not be optimal. A more sophisticated, truly structure-aware approach might distinguish between rare user-code functions and rare library functions, applying different generalization rules.
* \textbf{Missing Differential Privacy Baseline:} This study does not include a comparison to a formal Differential Privacy (DP) baseline. DP is the gold standard for privacy guarantees, and comparing our pragmatic approaches against a well-calibrated DP mechanism is a critical next step for benchmarking their position on the privacy-utility frontier.
* \textbf{Limited Threat Model:} Our privacy evaluation is based on a single, specific threat model: a project-linkage attack. While relevant, this does not encompass all possible privacy risks. Other attacks, such as attribute inference (e.g., inferring library versions or user configurations) or membership inference, were not evaluated.

\section{Conclusion}

\label{sec:conclusion}

This paper investigated the critical tradeoff between privacy and utility in the aggregation of software error traces. We proposed and implemented an evaluation framework to compare structured anonymization techniques, quantifying their impact using task-specific metrics for utility (Matching Accuracy) and re-identification risk (Leakage Risk). Our empirical study compared a pseudonymization approach (token hashing) with a generalization approach (a k-anonymity variant) across synthetic datasets of varying heterogeneity.

Our preliminary findings, based on a single experimental run, indicate that the choice of anonymization strategy should be informed by the nature of the data being protected. In a simulated high-diversity environment, representative of a large-scale error aggregator, the generalization method provided a better balance, achieving a Matching Accuracy of 0.6492 and a Leakage Risk of 0.0847. This performance was superior to token hashing in both utility and privacy in that setting. This suggests that for heterogeneous data, generalization's strategy of reducing identifier cardinality is more robust than pseudonymization's strategy of preserving it.

The main contribution of this work is the empirical comparison of these two anonymization philosophies in the context of error trace analysis. Our results provide initial, actionable evidence that developers of telemetry systems can and should move beyond collecting raw data, and that generalization-based techniques are a promising direction for heterogeneous environments.

Future work must proceed in three key directions. First, the experiments must be repeated with multiple random seeds to establish statistical significance and build confidence in the results. Second, the methods should be validated on a large, real-world dataset of error traces to ensure the findings generalize beyond our synthetic data. Finally, these pragmatic techniques must be benchmarked against a strong baseline with formal guarantees, such as a well-calibrated ε-Differential Privacy mechanism, to fully map the landscape of possible solutions.

\bibliographystyle{plainnat}
\bibliography{references}

\end{document}
